\subsection{Multiplexadores}

Os multiplexadores (MUX) atuam seleção do caminho de dados dentro do sistema.

No contexto deste projeto, os MUXs são utilizados para selecionar entre as duas possíveis saídas da lógica de transição de estados: incremento ou decremento. Dessa forma, a atualização do estado não é fixa, mas depende de um paramtro $\sigma$, controlado externamente, que define a operação a ser executada a cada ciclo de clock.

Formalmente, o multiplexador implementa a função $\delta(Q_t; \sigma)$:

\[
\delta(Q_t; sigma_t) = 
\begin{cases}
\delta_{inc}(Q_t), & \text{se } \sigma_t = 0 \\
\delta_{dec}(Q_t), & \text{se } \sigma_t = 1
\end{cases}
\]
